\makeabstract
    {
    Microwave Rotational Spectroscopy: Determining the Structure of 3-Chloro-2,3,3-Trifluoropropene Monomer and its Complexes with Argon and Acetylene
    }
    {
    \vspace{-20pt}
    Amaia C. Williams\textsuperscript{1}, Sophia Zhuang\textsuperscript{1}, Helen O. Leung\textsuperscript{1}, Mark D. Marshall\textsuperscript{1}
    }
    {
    \textsuperscript{1} Department of Chemistry, Amherst College, Amherst, MA
    }
    {
    Although weaker than chemical bonds, intermolecular interactions are essential for human life. The Leung and Marshall lab examines such interactions through exploring the configurations of halogenated ethenes and propenes. As a part of the lab{\textquoteright}s study of halopropenes, we seek to determine the structure of 3-chloro-2,3,3-trifluoropropene (3Cl233TFP).
    }
    {
    Gaussian16 was used to generate predicted structures of the monomer and its isotopologues. From the theoretical structure, rotational constants were obtained to generate a predicted spectrum. An experimental spectrum was obtained through a broadband, chirped-pulse Fourier transform microwave spectrometer. Comparing the two spectra in PGOPHER, the predicted rotational constants were adjusted to match experimental data. Experimental rotational constants for the most abundant species (and its isotopologues) were used to determine an experimental structure through Kraitchman coordinates and STRFIT calculations. Similar methods were used for 3Cl233TFP{\textquoteright}s complex with argon and acetylene. For the acetylene complex, basis set superposition error (BSSE) and zero point energy (ZPE) corrections were applied to determine the lowest energy configuration. 
    }
    {
    The potential energy scan of the 3Cl233TFP monomer yielded two unique rotamers, from which the lowest energy configuration was determined. The spectra of both rotamers and their Cl-37 isotopologues were assigned to obtain rotational constants; C-13 isotopologues were only assigned for the lowest energy rotamer. An experimental structure for the monomer was obtained. An energy scan of argon around 3Cl233TFP{\textquoteright}s rotamers resulted in 3 distinct minima for the lower energy configuration and 1 distinct minima for the higher energy configuration. Experimental rotational constants were only determined for the lowest energy Argon complex (r = 4.0132, \ensuremath{\theta} = 134.7254, \ensuremath{\Phi} = 65.6724) and its Cl-37 isotopologue. An experimental structure of the Ar-3Cl233TFP complex was determined. An energy scan of acetylene around the lowest energy rotamer resulted in 3 distinct minima. BSSE and ZPE corrections were applied to identify the absolute minima at (r = 4.1781, \ensuremath{\theta} = 71.8738, \ensuremath{\Phi} = 234.7507). 
    }
    {
    Theoretical predictions of the 3Cl233TFP monomer and its complexes were adequate, as experimental and predicted structures yielded similar results. Further explorations of the HCCH-3Cl233TFP complex are ongoing to determine the experimental structure of the HCHH-3Cl233TFP structure.
    }

    \newpage
    